% Baccalauréat malgache, série D, 2016 : énoncé.
\TCDBacAttribution{D}{2016}
\providecommand{\Plan}{(\mathscr{P})}
\providecommand{\Ouv}{(O ; \vec{u}, \vec{v})}
\providecommand{\cm}{~\mathrm{cm}}
\providecommand{\Courbe}{(\mathscr{C})}
\providecommand{\CourbeP}{(\mathscr{C}')}
\providecommand{\D}{(\mathscr{D})}
\providecommand{\dx}{~\mathrm{d}x}

\noindent \underline{\textsc{Exercice I}} (5 points) \newline
      Le plan complexe $\Plan$ est muni d'un rep\`ere orthonorm\'e direct $\mathcal{R} = \Ouv$ d'unit\'e $2\cm$.
      \begin{enumerate}
            \item $P$ est le polyn\^ome \`a variable complexe $z$ d\'efini par : $P(z) = z^3 - 2z^2 + (2+i) z - 1 - i$.
                  \begin{enumerate}
                        \item Calculer $P(1)$ et $(1-2i)^2$.
                        \item R\'esoudre dans $\mathbb{C}$ l'\'equation $P(z) = 0$.
                  \end{enumerate}
            \item On consid\`ere les points $A$, $B$ et $C$ tels que $z_A = 1$, $z_B = 1-i$ et $z_C = i$.
                  Ecrire le nombre complexe $a = \frac{z_B}{z_C}$ sous forme trigonom\'etrique et en d\'eduire $a^4$.
            \item On consid\`ere la similitude plane directe $S$ qui transforme $M$ en $M'$ telle que : 
                  $z' = (1-i) z + 2i$ ; $M'$ d'affixe $z'$ et $M$ d'affixe $z$.
                  \begin{enumerate}
                        \item D\'eterminer les \'el\'ements caract\'eristiques de $S$.
                        \item Donner la nature et les \'el\'ements caract\'eristiques de $S^4 = S \circ S \circ S \circ S$.
                        \item Construire dans le m\^eme rep\`ere le triangle $ABC$ et $A'B'C'$ son image par $S^4$.
                  \end{enumerate}
      \end{enumerate}
%%%%%%%%%%%%%%%%%%%%%%%%%%%%%%%%%%%%%%%%%%%%%%%%%%%%%%%%%%%%%%%%%%%%%%%%%%%%%%%%%%%%%%%%%%%%%%%
%%%%%%%%%%%%%%%%%%%                   %%%%%%%%%%%%%%%%%%%                   %%%%%%%%%%%%%%%%%%%
%%%%%%%%%%%%%%%%%%%%%%%%%%%%%%%%%%%%%%%%%%%%%%%%%%%%%%%%%%%%%%%%%%%%%%%%%%%%%%%%%%%%%%%%%%%%%%%
\noindent \underline{\textsc{Exercice II}} (5 points) 
      \begin{enumerate}
            \item Les faces d'un d\'e cubique truqu\'e sont num\'erot\'ees : 1, 2, 2, 3, 4, 6.
                  On lance une fois ce d\'e. 
                  Chaque face a la m\^eme probabilit\'e d'apparition.
                  On note $P_i$ la probabilit\'e d'apparition de la face portant le num\'ero $i$. 
                  Calculer $P_1$, $P_2$, $P_3$, $P_4$, $P_6$.
            \item On lance deux fois de suite ce d\'e. 
                  Chaque face a toujours la m\^eme probabilit\'e d'apparition.
                  \begin{enumerate}
                        \item Calculer la probabilit\'e de l'\'ev\`enement :
                              $A$ : << La somme des deux num\'eros apparus sur la face sup\'erieure est \'egale \`a 4
                                    lors de deux lancements de ce d\'e >>.
                        \item On note par $X$ la variable al\'eatoire d\'efinie par la somme des num\'eros affich\'es 
                              lors de deux lancements de ce d\'e.
                              \begin{itemize}
                                    \item Donner la loi de probabilit\'e de $X$.
                                    \item Calculer l'esp\'erance math\'ematique $E(X)$.
                              \end{itemize}
                        \item On lance trois fois de suite et d'une mani\`ere ind\'ependante ce d\'e.
                              On note par $Y$ la variable al\'eatoire \'egale au nombre d'apparitions de la face portant
                              le num\'ero 2 lors de ces trois lancers.
                              \begin{enumerate}
                                    \item Donner la loi de probabilit\'e de $Y$.
                                    \item Calculer la variance $V(Y)$.
                              \end{enumerate}
                  \end{enumerate}
            \textbf{NB : Mettre les r\'esultats sous forme de fraction irr\'eductible.}
      \end{enumerate}
%%%%%%%%%%%%%%%%%%%%%%%%%%%%%%%%%%%%%%%%%%%%%%%%%%%%%%%%%%%%%%%%%%%%%%%%%%%%%%%%%%%%%%%%%%%%%%%
%%%%%%%%%%%%%%%%%%%                   %%%%%%%%%%%%%%%%%%%                   %%%%%%%%%%%%%%%%%%%
%%%%%%%%%%%%%%%%%%%%%%%%%%%%%%%%%%%%%%%%%%%%%%%%%%%%%%%%%%%%%%%%%%%%%%%%%%%%%%%%%%%%%%%%%%%%%%%
\noindent \underline{\textsc{Probleme}} \newline
      \textbf{NB : Les parties A et B sont ind\'ependantes.}
      \begin{description}
            \item[A. ]
                  Soit $f$ la fonction d\'efinie sur $\mathbb{R}$ par : $f(x) = x \expo{-x} + x - 1$.
                  On note par $\Courbe$ sa courbe repr\'esentative dans un rep\`ere orthonorm\'e \Oij{} d'unit\'e $2\cm$.
                  \begin{enumerate}
                        \item Soit $g$ la fonction d\'efinie sur $\mathbb{R}$ par : $g(x) = (1-x) \expo{-x} + 1$.
                              \begin{enumerate}
                                    \item Etudier le sens de variation de $g$.
                                    \item En d\'eduire le signe de $g(x)$.
                              \end{enumerate}
                        \item Etudier les variations de $f$.
                        \item \begin{enumerate}
                                    \item Etudier la branche infinie de $\Courbe$ au voisinage de $-\infty$.
                                    \item D\'emontrer que $\Courbe$ admet une asymptote oblique dont on pr\'ecisera
                                          l'\'equation au voisinage de $+\infty$ qu'on nommera $\D$.
                                          On pr\'ecisera la position de $\Courbe$ par rapport \`a $\D$.
                              \end{enumerate}
                        \item D\'emontrer qu'il existe un point unique $A$ de $\Courbe$ o\`u la tangente $(T)$ est 
                              parall\`ele \`a $\D$.
                        \item Tracer dans le m\^eme rep\`ere $\D$, $(T)$ et $\Courbe$.
                        \item Calculer en $\cm^2$, l'aire g\'eom\'etrique du domaine plan limit\'e par la courbe $\Courbe$,
                              la droite $\D$ et les droites d'\'equations $x = 0$ et $x = 1$.
                        \item \begin{enumerate}
                                    \item D\'emontrer que $f$ r\'ealise une bijection de $\mathbb{R}$ sur un intervalle
                                          $J$ que l'on d\'eterminera.
                                    \item Calculer $f(1)$ et $(f^{-1})'\left( \frac{1}{e} \right)$ o\`u $f^{-1}$ est 
                                          la fonction r\'eciproque de $f$.
                                    \item Tracer dans le m\^eme rep\`ere que $\Courbe$ la courbe $\CourbeP$ de $f^{-1}$.
                                          (On donne : $e = 2,7$, $\frac{1}{e} = 0,4$).
                              \end{enumerate}
                  \end{enumerate}
            \item[B. ]
                  Soit $(U_n)_{n\in\mathbb{N}}$ une suite d\'efinie par $U_0 = 1$ et $U_n = 1 - \int_{0}^{n} \expo{-x} \dx$.
                  \begin{enumerate}
                        \item Calculer $U_n$ en fonction de $n$.
                        \item D\'emontrer que $(U_n)_{n\in\mathbb{N}}$ 
                              est une suite g\'eom\'etrique dont on pr\'ecisera la raison.
                  \end{enumerate}
      \end{description}










